% Section 3 - RAII
% Roberto Masocco <roberto.masocco@uniroma2.it>
% April 12, 2026

% ### RAII ###
\section{RAII}
\graphicspath{{figs/section3/}}

% --- RAII ---
\begin{frame}{RAII}{The resource problem}
	\justifying
	A \textbg{resource} is anything that must be \textbg{acquired before use} and \textbg{released after use}:
	\begin{itemize}
		\item \textbg{memory} allocated on the heap (\texttt{malloc}/\texttt{free} in C, \texttt{new}/\texttt{delete} in C++);
		\item \textbg{mutex locks} acquired with \texttt{pthread\_mutex\_lock};
		\item \textbg{file descriptors} opened with \texttt{fopen}/\texttt{open};
		\item \textbg{network sockets}, hardware handles, GPU buffers...
	\end{itemize}
	\bigskip
	In C, the programmer is responsible for releasing every resource on \textbg{every possible exit path} from a function — including error paths.\\
	\bigskip
	\begin{alertblock}{The classic bug}
		\justifying
		An early \texttt{return}, a \texttt{goto}, or an unhandled error condition leaves a \textbr{resource unreleased}.\\
		In a long-running system this causes \textbr{resource leaks} that accumulate over time.
	\end{alertblock}
\end{frame}
\begin{frame}{RAII}{Resource Acquisition Is Initialization}
	\begin{block}{Definition of RAII}
		\justifying
		\textbf{Resource Acquisition Is Initialization (RAII)} is a programming idiom in which a resource is \textbf{acquired in the constructor} of an object and \textbf{released in the destructor}.
	\end{block}
	\justifying
	\bigskip
	Because the destructor is \textbg{guaranteed to run} when an object goes out of scope, resource release becomes \textbg{automatic and unconditional}:
	\begin{itemize}
		\item no matter how many \texttt{return} statements exist in a function;
		\item no matter which branch of an \texttt{if}/\texttt{switch} is taken;
		\item even if an \textbg{exception} is thrown.
	\end{itemize}
	\bigskip
	\begin{block}{}
		\centering
		\textbf{RAII is the most important C++ programming idiom.}\\
		Every major abstraction in the standard library is built on it.
	\end{block}
\end{frame}
\begin{frame}{RAII}{RAII in practice}
	\justifying
	The idiom is not a language feature — it is a \textbg{design pattern} that the language makes reliable through its object lifetime rules.\\
	\smallskip
	Some examples of objects that use this idiom in C++, from your OS course:
	\begin{itemize}
		\item a \textbg{mutex wrapper} that locks on construction and unlocks on destruction;
		\item a \textbg{file handle} class that opens on construction and closes on destruction;
		\item a \textbg{memory buffer} class that allocates on construction and deallocates on destruction.
	\end{itemize}
	\smallskip
	In C++ the standard library provides \textbg{ready-made RAII wrappers} for all common resources:
	\begin{itemize}
		\item \texttt{std::unique\_ptr}, \texttt{std::shared\_ptr} for heap memory;
		\item \texttt{std::lock\_guard}, \texttt{std::unique\_lock} for mutexes;
		\item \texttt{std::fstream} for files.
	\end{itemize}
\end{frame}
